\chapter{Faire un marché}\label{ch:iv2:making-a-market}

\begin{ivtrains}
Donner un cours acheteur et un cours vendeur sur n'importe quoi, d'un coup, et les déplacer pour la bonne raison~: la valeur, la distribution, qui se trouve de l'autre côté, ce que leur transaction vient d'apprendre et ce que l'on détient déjà. Dans chaque question, le cours acheteur est le prix auquel on achète et le cours vendeur celui auquel on vend. L'énoncé fixe la largeur de la cotation et ce que fait l'autre côté, si bien que chaque question admet une seule réponse vérifiable. Les premières questions prennent quelques minutes, et les dernières sont des dérivations.
\end{ivtrains}

\begin{question}[$\star$ \iqroles{QT, Trader, QR} \iqfirm{teneur de marché, hedge fund, crypto} \iqlevel{junior}]\label{iq:iv2:making-a-market:1}
Je lance trois dés équilibrés, et un contrat paie leur somme en dollars. Cotez un cours acheteur et un cours vendeur à exactement 1 d'écart, par pas de 0.5. Un client qui n'a rien vu achète ou vend alors un lot à vos prix, le côté étant choisi à pile ou face. Que cotez-vous, et quel gain espérez-vous~? Une cotation à 11 en cours acheteur et 12 en cours vendeur rapporterait-elle moins en moyenne~?
\end{question}

\begin{question}[$\star$ \iqroles{QT, Trader, QR} \iqfirm{teneur de marché, hedge fund} \iqlevel{junior}]\label{iq:iv2:making-a-market:2}
Je tire une carte d'un jeu de 52 cartes mélangé et je la regarde. Un contrat paie son rang, de l'as valant 1 au roi valant 13. Cotez un cours acheteur et un cours vendeur à exactement 2 d'écart, par pas de 0.5. Je vous vends un lot à votre cours acheteur si le rang lui est inférieur, j'achète un lot à votre cours vendeur s'il est supérieur, et sinon je passe. Quelle cotation perd le moins en moyenne, et combien perd-elle~? Si vous pouvez choisir une largeur entière quelconque, quelle est la plus étroite qui perd moins de 1~?
\end{question}

\begin{question}[$\star$ \iqroles{QT, Trader, QR} \iqfirm{teneur de marché, hedge fund, crypto} \iqlevel{junior}]\label{iq:iv2:making-a-market:3}
Un dé équilibré est lancé, et un contrat paie le carré du résultat. Cotez un cours acheteur et un cours vendeur à exactement 2 d'écart, par pas de 0.5. Un client qui n'a pas vu le dé achète ou vend un lot, le côté étant choisi à pile ou face. Que cotez-vous~?
\end{question}

\begin{question}[$\star$ \iqroles{QT, Trader} \iqfirm{teneur de marché, hedge fund} \iqlevel{junior}]\label{iq:iv2:making-a-market:4}
Dans un jeu de 52 cartes, comptez l'as pour 1, les cartes numérotées à leur valeur faciale, le valet pour 11, la dame pour 12 et le roi pour 13. Dix cartes ont été retournées face visible, et elles totalisent 100. Un contrat paie la valeur de la carte suivante, et un autre paie la somme des 42 cartes restées face cachée. Vous cotez chacun avec une largeur de 1, par pas de 0.25. Votre contrepartie voit les mêmes dix cartes. Elle achète à votre cours vendeur si celui-ci est inférieur à l'espérance du contrat, vend à votre cours acheteur s'il est supérieur, et sinon passe. Quelles sont vos deux cotations~? Quel contrat vous fait-il courir le plus de risque~?
\end{question}

\begin{question}[$\star$ \iqroles{QT, Trader, QR} \iqfirm{teneur de marché, crypto} \iqlevel{junior}]\label{iq:iv2:making-a-market:5}
Un dé va être lancé. Le contrat $A$ paie 1~dollar si le résultat est 4, 5 ou 6, et le contrat $B$ paie 1~dollar s'il vaut 1, 2 ou 3. Vous affichez $A$ à 0.55 en cours acheteur et 0.60 en cours vendeur, et $B$ à 0.50 en cours acheteur et 0.52 en cours vendeur. Votre contrepartie prend toute combinaison de transactions d'un lot qui garantit un gain certain. Que fait-elle, et combien gagne-t-elle~? Que doivent satisfaire vos quatre prix pour qu'aucune combinaison de ce type n'existe~?
\end{question}

\begin{question}[$\star$ \iqroles{QT, Trader} \iqfirm{teneur de marché, hedge fund, crypto} \iqlevel{junior}]\label{iq:iv2:making-a-market:6}
Deux dés équilibrés sont lancés l'un après l'autre, et un contrat paie le plus élevé des deux. Cotez un cours acheteur et un cours vendeur à exactement 1 d'écart, par pas de 0.25, avant que l'un ou l'autre dé ne soit lancé. Un client qui n'a rien vu achète ou vend un lot à pile ou face. Puis le premier dé montre un 3, et vous cotez à nouveau selon les mêmes modalités. Quelles sont vos deux cotations~? Quelle aurait été votre seconde cotation si le premier dé avait montré un 6~?
\end{question}

\begin{question}[$\star$ \iqroles{QT, Trader, QR} \iqfirm{teneur de marché, hedge fund} \iqlevel{junior}]\label{iq:iv2:making-a-market:7}
Faites un marché sur le nombre de nombres premiers inférieurs à 1\,000, sans calculatrice. Un contrat paie ce nombre en dollars. Votre cours acheteur et votre cours vendeur doivent être à exactement 20 d'écart, en nombres entiers. Je connais la réponse~: je vous vends un lot à votre cours acheteur si le vrai nombre est inférieur à celui-ci, j'achète un lot à votre cours vendeur si le vrai nombre est supérieur à celui-ci, et sinon je passe. Que cotez-vous, et comment l'avez-vous obtenu~?
\end{question}

\begin{question}[$\star\star$ \iqroles{QT, Trader, QR} \iqfirm{teneur de marché, hedge fund, crypto} \iqlevel{junior}]\label{iq:iv2:making-a-market:8}
Vous lancez un dé équilibré jusqu'à l'apparition du premier 6, et un contrat paie le nombre de lancers en dollars. Cotez un cours acheteur et un cours vendeur à exactement 2 d'écart, en nombres entiers. Un client qui ne sait rien achète ou vend un lot à pile ou face. Que cotez-vous~? Comparez les deux issues possibles de la transaction~: dans chacune, quel gain espérez-vous, quelle est la probabilité de perdre de l'argent et combien pouvez-vous perdre~?
\end{question}

\begin{question}[$\star\star$ \iqroles{QT, Trader, QR} \iqfirm{teneur de marché, hedge fund} \iqlevel{junior}]\label{iq:iv2:making-a-market:9}
Un contrat paie le nombre de cartes rouges parmi les 10 cartes du dessus d'un jeu de 52 cartes mélangé. Cotez un cours acheteur et un cours vendeur à exactement 2 d'écart, par pas de 0.5. Un client qui n'a pas vu les cartes achète ou vend un lot à pile ou face. Que cotez-vous, et quelle est la probabilité que le nombre final tombe en dehors de votre cotation~? La première carte est ensuite retournée, et c'est une carte rouge. Où cotez-vous maintenant, aux mêmes conditions~?
\end{question}

\begin{question}[$\star\star$ \iqroles{QT, Trader, QR} \iqfirm{teneur de marché, hedge fund} \iqlevel{junior, confirmé}]\label{iq:iv2:making-a-market:10}
Vous pensez qu'une quantité inconnue $X$ suit une loi normale de moyenne $\mu$ et d'écart type $\sigma$. Vous cotez un cours acheteur de $\mu - h$ et un cours vendeur de $\mu + h$, et un contrat se règle en $X$. Votre contrepartie connaît $X$~: elle vous vend un lot à votre cours acheteur si $X$ lui est inférieur, achète un lot à votre cours vendeur si $X$ lui est supérieur, et sinon passe. Calculez votre perte espérée en fonction de $h$. Combien perdez-vous si votre cotation est l'intervalle de confiance à 90~\%~? Quelle largeur maintient la perte sous 5~\% de $\sigma$~?
\end{question}

\begin{question}[$\star\star$ \iqroles{QT, Trader, QR} \iqfirm{teneur de marché, hedge fund} \iqlevel{junior, confirmé}]\label{iq:iv2:making-a-market:11}
Un contrat paie, en dollars, le nombre d'accordeurs de piano d'une ville de cinq millions d'habitants. Votre raisonnement approximatif donne une estimation de 200, et vous laisse incertain d'un facteur 2 environ. Considérez que le logarithme népérien du nombre suit une loi normale, de moyenne $\ln 200$ et d'écart type $\ln 2$. Vous pouvez coter n'importe quels prix, et un client qui ne sait rien achète ou vend un lot à pile ou face. Autour de quel nombre centrez-vous votre cotation, et pourquoi pas 200~? Quelle est la probabilité que le vrai nombre dépasse votre centre~? Quel intervalle contient le vrai nombre avec une probabilité de 90~\%~?
\end{question}

\begin{question}[$\star\star$ \iqroles{QT, Trader, QR} \iqfirm{teneur de marché, hedge fund, crypto} \iqlevel{junior}]\label{iq:iv2:making-a-market:12}
Un contrat vaut soit 80, soit 120, avec une probabilité égale. Vous cotez un cours acheteur et un cours vendeur, à des prix quelconques, et un trader arrive et échange un lot. Avec une probabilité de 30~\%, il connaît la vraie valeur et achète si elle vaut 120, vend si elle vaut 80. Sinon, il ne sait rien et achète ou vend à pile ou face. À quels cours acheteur et vendeur êtes-vous exactement à l'équilibre en moyenne, quel que soit le côté choisi~?
\end{question}

\begin{question}[$\star\star$ \iqroles{QT, Trader, QR} \iqfirm{teneur de marché, hedge fund, crypto} \iqlevel{junior, confirmé}]\label{iq:iv2:making-a-market:13}
Un contrat vaut soit 0, soit 100, avec une probabilité égale. Les traders arrivent un à un, et chacun est une nouvelle personne. Chacun échange un lot~: avec une probabilité de 20~\%, il connaît la valeur et achète si elle vaut 100, vend si elle vaut 0. Sinon, il achète ou vend à pile ou face. Avant chaque trader, vous cotez le cours acheteur et le cours vendeur auxquels vous êtes à l'équilibre en moyenne, à des prix quelconques. Les quatre premières transactions sont~: achat, achat, vente, achat. Quelle est votre cotation avant chacun des cinq premiers traders~? Si les mêmes quatre transactions étaient arrivées dans un autre ordre, votre cotation finale changerait-elle~?
\end{question}

\begin{question}[$\star\star$ \iqroles{QT, Trader, QR} \iqfirm{teneur de marché, hedge fund} \iqlevel{junior}]\label{iq:iv2:making-a-market:14}
Un dé équilibré a été lancé hors de votre vue, et un contrat paie sa valeur. Vous cotez 3 en cours acheteur et 4 en cours vendeur. Votre contrepartie a regardé le dé avec une probabilité d'un demi. Si elle l'a regardé, elle vous vend un lot à votre cours acheteur si la valeur est strictement inférieure à celui-ci, achète un lot à votre cours vendeur si elle est strictement supérieure à celui-ci, et sinon passe. Sinon, elle achète ou vend à pile ou face. Vous ne savez pas laquelle des deux situations s'est réalisée. Elle vous vend. Que vaut le contrat maintenant~? Où cotez-vous ensuite, avec une largeur de 1, centrée sur votre nouvelle valeur, par pas de 0.1~? Que vaudrait le contrat si elle avait passé à la place~?
\end{question}

\begin{question}[$\star\star$ \iqroles{QT, Trader} \iqfirm{teneur de marché, hedge fund} \iqlevel{junior, confirmé}]\label{iq:iv2:making-a-market:15}
Un contrat vaut 90 ou 110, avec une probabilité égale. Un trader arrive. Avec une probabilité de 10~\%, il connaît la valeur et achète si elle vaut 110, vend si elle vaut 90. Sinon, il achète ou vend à pile ou face. Un trader qui connaît la valeur échange 100 lots avec une probabilité de 80~\% et 1 lot sinon. Un trader qui ne la connaît pas échange 100 lots avec une probabilité de 5~\% et 1 lot sinon. Vous voyez la taille et le côté, et vous pouvez annoncer un prix différent pour chaque taille. À quel cours vendeur êtes-vous à l'équilibre en moyenne pour un acheteur d'un lot, et pour un acheteur de 100 lots~? Si vous devez afficher un seul cours vendeur pour les deux tailles, quel est-il~?
\end{question}

\begin{question}[$\star\star$ \iqroles{QT, Trader, QR} \iqfirm{teneur de marché, hedge fund, crypto} \iqlevel{junior, confirmé}]\label{iq:iv2:making-a-market:16}
La juste valeur d'un contrat est 100 et ne changera pas avant son règlement à 100. Vous êtes long de 40 lots, et votre score final est votre profit moins $0.01$ fois le carré de votre position finale. Un client peut acheter un lot à votre cours vendeur ou en vendre un à votre cours acheteur, ou bien ne rien faire. Quel est le cours vendeur le plus bas et le cours acheteur le plus haut auxquels une transaction ne ferait pas baisser votre score~? Votre cotation doit avoir une largeur exactement de 2, centrée entre ces deux prix~: qu'affichez-vous~? Qu'affichez-vous si vous êtes short de 40~? De combien votre centre se déplace-t-il après chaque transaction~?
\end{question}

\begin{question}[$\star\star$ \iqroles{QT, Trader} \iqfirm{teneur de marché, crypto} \iqlevel{junior}]\label{iq:iv2:making-a-market:17}
La juste valeur d'un contrat est fixée à 100, et vous cotez toujours 99 en cours acheteur et 101 en cours vendeur. À chaque manche, exactement un client arrive et achète ou vend un lot à pile ou face. Votre position doit rester entre $-3$ et $+3$ lots~: à $+3$, vous baissez votre cours acheteur, et à $-3$, vous baissez votre cours vendeur, de sorte qu'un client de ce côté s'en va sans transaction. Vous partez à plat, et toute position finale est évaluée à 100. Quel est votre profit espéré sur 10 manches~? Quel est-il par manche à long terme~? Quel est-il par manche à long terme avec une limite de $\pm L$ lots~?
\end{question}

\begin{question}[$\star\star$ \iqroles{QT, Trader, QD} \iqfirm{teneur de marché, banque, crypto} \iqlevel{junior}]\label{iq:iv2:making-a-market:18}
Un contrat $S$ paie la somme de trois dés équilibrés. Un autre dealer cote un contrat sur chaque dé seul à 3.4 en cours acheteur et 3.6 en cours vendeur, en taille quelconque, et continuera à le faire. Vous devez afficher une cotation sur $S$. Votre contrepartie prend toutes les transactions qui garantissent un gain certain contre vous et contre ce dealer. Pour quels cours acheteur et quels cours vendeur peut-elle le faire~? Un client distinct, qui ne sait rien, achète alors un lot de $S$ auprès de vous à votre cours vendeur de 10.7. Que gagnez-vous si vous couvrez immédiatement avec les dés seuls, et que gagnez-vous si vous ne couvrez pas~?
\end{question}

\begin{question}[$\star\star$ \iqroles{QT, Trader, QR} \iqfirm{teneur de marché, hedge fund} \iqlevel{junior, confirmé}]\label{iq:iv2:making-a-market:19}
Trois dés équilibrés $X$, $Y$ et $Z$ sont lancés hors de votre vue. Le contrat $A$ paie $X + Y$ et le contrat $B$ paie $X + Z$, et vous cotez chacun à 6.5 en cours acheteur et 7.5 en cours vendeur. Votre contrepartie connaît la somme $X + Y$ et rien d'autre. Elle achète un lot de $A$ à votre cours vendeur si cette somme est supérieure à celui-ci, vend à votre cours acheteur si elle est inférieure, et sinon passe. Elle vous achète $A$. Que vaut chaque contrat maintenant~? Compte tenu de ce que vous savez désormais, combien de lots de $B$, fractions autorisées, devez-vous acheter pour minimiser la variance de votre position, et quelle part de celle-ci cela supprime-t-il~?
\end{question}

\begin{question}[$\star\star$ \iqroles{QT, Trader, QD, LLE} \iqfirm{teneur de marché, crypto} \iqlevel{junior, confirmé}]\label{iq:iv2:making-a-market:20}
Vous cotez un contrat à sa juste valeur plus et moins $h$. Entre une information et la suivante, 10 clients non informés échangent chacun un lot avec vous, et chacun vous rapporte $h$. Ensuite, l'information suivante déplace la juste valeur de 2, à la hausse ou à la baisse. Avec une probabilité de 80~\%, un trader rapide frappe votre ancienne cotation du côté devenu faux, pour un lot, avant que vous ne puissiez la déplacer. Sinon, vous réagissez à temps. Quel $h$ vous met à l'équilibre~? Quel est-il si un système plus rapide que le vôtre ramène cette probabilité de 80~\% à 50~\%, et s'il n'y a que 5 clients~?
\end{question}

\begin{question}[$\star\star$ \iqroles{QT, Trader, QR} \iqfirm{teneur de marché, crypto} \iqlevel{junior, confirmé}]\label{iq:iv2:making-a-market:21}
Un contrat vaut 98 ou 102, avec une probabilité égale. Chaque trader échange un lot~: 10~\% d'entre eux connaissent la valeur et achètent si elle vaut 102 et vendent si elle vaut 98, et les autres achètent ou vendent à pile ou face. Un concurrent affiche 99 en cours acheteur et 101 en cours vendeur. Vous pouvez vous joindre aux mêmes prix, auquel cas chaque transaction est répartie à parts égales entre vous et lui. Ou bien vous affichez $100 - h$ en cours acheteur et $100 + h$ en cours vendeur, pour un certain $h$ inférieur à 1, par pas de 0.1, et vous prenez chaque transaction. Quel est votre gain par trader dans chaque cas, et quel $h$ est le meilleur si le rival reste en place~? Où les cotations finissent-elles si le rival bouge aussi, un pas à la fois~?
\end{question}

\begin{question}[$\star\star$ \iqroles{QT, Trader, QR} \iqfirm{teneur de marché, hedge fund} \iqlevel{junior, confirmé}]\label{iq:iv2:making-a-market:22}
Un contrat paie la somme de trois dés équilibrés, et vous cotez $10.5 - h$ en cours acheteur et $10.5 + h$ en cours vendeur. Avec une probabilité d'un demi, votre contrepartie a regardé un dé. Elle calcule alors l'espérance de la somme, vend un lot si celle-ci est inférieure à votre cours acheteur, achète un lot si elle est supérieure à votre cours vendeur, et sinon passe. Sinon, elle n'a rien vu et achète ou vend à pile ou face. Quel $h$ vous met à l'équilibre en moyenne~? Quel est-il si elle regarde avec une probabilité d'un quart, ou de trois quarts~?
\end{question}

\begin{question}[$\star\star\star$ \iqroles{QT, Trader, QR} \iqfirm{teneur de marché, hedge fund, crypto} \iqlevel{junior, confirmé}]\label{iq:iv2:making-a-market:23}
Un contrat vaut 0 ou 100, avec une probabilité égale. Les traders arrivent un à un, et chacun est une nouvelle personne. La moitié d'entre eux connaissent la valeur, et achètent si elle vaut 100 et vendent si elle vaut 0. L'autre moitié achète ou vend à pile ou face. Avant chaque trader, vous cotez le cours acheteur et le cours vendeur auxquels vous êtes à l'équilibre en moyenne, à des prix quelconques. Jusqu'ici, tous les traders ont vendu. Quelle est votre cotation avant les cinq premiers traders~? Quand votre largeur passe-t-elle pour la première fois sous 5~? Comment la largeur évolue-t-elle à mesure que les ventes se poursuivent, et pourquoi~? Quelle est la probabilité d'une série de quatre ventes si la valeur vaut réellement 100~?
\end{question}

\begin{question}[$\star\star\star$ \iqroles{QT, QR} \iqfirm{teneur de marché, hedge fund} \iqlevel{junior, confirmé}]\label{iq:iv2:making-a-market:24}
La valeur $V$ d'un contrat est uniforme sur 0 à 100. Un trader arrive et échange un lot. Avec une probabilité $\alpha$, il connaît $V$~: il achète à votre cours vendeur si $V$ est supérieure à celui-ci, vend à votre cours acheteur si $V$ est inférieure à celui-ci, et sinon passe. Avec une probabilité $1 - \alpha$, il ne sait rien et achète ou vend à pile ou face. Déterminez le cours vendeur auquel un achat vous met à l'équilibre en moyenne, et le cours acheteur auquel une vente le fait, en fonction de $\alpha$. Évaluez-les pour $\alpha = 0.5$. Que se passe-t-il quand $\alpha$ tend vers 1~?
\end{question}

\begin{question}[$\star\star\star$ \iqroles{QT, Trader, QR} \iqfirm{teneur de marché, hedge fund} \iqlevel{confirmé}]\label{iq:iv2:making-a-market:25}
La valeur $V$ d'un contrat est uniforme sur 0 à 100, et vous cotez $50 - h$ en cours acheteur et $50 + h$ en cours vendeur. Un trader arrive. Avec une probabilité de 20~\%, il connaît $V$~: il achète un lot à votre cours vendeur si $V$ est supérieure à celui-ci, vend un lot à votre cours acheteur si $V$ est inférieure à celui-ci, et sinon passe. Avec une probabilité de 80~\%, il ne sait rien~: il choisit un côté à pile ou face et échange un lot avec une probabilité $1 - h/20$, et pas du tout si $h$ vaut 20 ou plus. Quel $h$ maximise votre profit espéré, et quelle est la valeur de ce profit~? Pour quel $h$ plus petit vous ne faites que l'équilibre~? Au-delà de quelle part de traders informés aucune cotation n'est-elle rentable~?
\end{question}

\begin{question}[$\star\star\star$ \iqroles{QT, Trader, QR} \iqfirm{teneur de marché, crypto} \iqlevel{junior, confirmé}]\label{iq:iv2:making-a-market:26}
La juste valeur d'un contrat est 100 et reste à ce niveau jusqu'à son règlement. Votre score est votre profit moins $0.01$ fois le carré de votre position finale. Un seul client arrive, voulant acheter ou vendre un lot à pile ou face. Un vendeur prend votre cours acheteur avec une probabilité $e^{-2\delta}$, où $\delta$ est la distance à laquelle votre cours acheteur se situe sous 100. Un acheteur prend votre cours vendeur avec une probabilité $e^{-2\delta}$, $\delta$ étant mesuré au-dessus de 100. Un cours acheteur supérieur à 100 ou un cours vendeur inférieur à 100 est pris à coup sûr. Quels cours acheteur et vendeur maximisent votre score espéré quand vous êtes à plat, et quand vous êtes long de 40~? À partir de quelle position votre cours vendeur atteint-il 100~?
\end{question}

\begin{question}[$\star\star\star$ \iqroles{QT, Trader, QR} \iqfirm{teneur de marché, hedge fund} \iqlevel{junior, confirmé}]\label{iq:iv2:making-a-market:27}
Un contrat vaut 90 ou 110, avec la même probabilité. Un trader arrive~: avec une probabilité de 30~\%, il connaît la valeur et achète si elle vaut 110, vend si elle vaut 90, et sinon il achète ou vend à pile ou face. Quels cours acheteur et vendeur à l'équilibre cotez-vous~? Maintenant vous voyez aussi un signal public, vu de tous, qui donne la vraie valeur correctement avec une probabilité de 80~\%. Il indique 110. Quelle cotation à l'équilibre affichez-vous maintenant~? Comparez les deux largeurs de la fourchette. Quelle largeur le signal laisse-t-il s'il n'est exact que dans 60~\% des cas~?
\end{question}

\begin{question}[$\star\star\star$ \iqroles{QT, Trader, QR} \iqfirm{teneur de marché, hedge fund, crypto} \iqlevel{junior, confirmé}]\label{iq:iv2:making-a-market:28}
Un contrat rapporte la somme de deux dés équilibrés. Vous ne voyez que le premier dé, et un dealer concurrent ne voit que le second. Chacun de vous propose de vendre un lot à son propre dé plus 3.5 plus une marge $m$, la même $m$ pour les deux, et vous connaissez tous deux cette règle. Un client achète toujours un lot au cours vendeur le plus bas, et si les cours vendeurs sont égaux, chacun de vous vend un demi-lot. Quel est votre profit espéré pour $m = 0.5$~? Pour quel $m$ êtes-vous à l'équilibre~? Avec trois dealers et trois dés, chacun voyant un dé et proposant son propre dé plus 7 plus $m$, les cours vendeurs les plus bas se partageant la vente à parts égales, pour quel $m$ êtes-vous à l'équilibre~?
\end{question}

\begin{question}[$\star\star\star$ \iqroles{QT, Trader, QR} \iqfirm{teneur de marché, hedge fund} \iqlevel{junior, confirmé}]\label{iq:iv2:making-a-market:29}
Un contrat vaut 0 ou 100, avec la même probabilité. Une personne échange avec vous à chaque manche, un lot à la fois. Avec une probabilité d'un demi, cette personne connaît la valeur et achète toujours si elle vaut 100, vend toujours si elle vaut 0. Sinon elle ne sait rien et choisit chaque côté à pile ou face. Avant chaque manche, vous cotez le cours acheteur et le cours vendeur auxquels vous êtes à l'équilibre en moyenne pour cette manche, quels que soient les prix. Que cotez-vous à la manche 1~? Et après qu'elle a acheté une fois, puis après qu'elle a acheté puis vendu~? Comparez avec le cas où une nouvelle personne arrive à chaque manche, dont la moitié connaît la valeur.
\end{question}

\begin{question}[$\star\star\star$ \iqroles{QT, Trader, QR} \iqfirm{teneur de marché, hedge fund} \iqlevel{confirmé}]\label{iq:iv2:making-a-market:30}
Dix cartes numérotées de 1 à 10 sont mélangées, et trois sont distribuées face cachée. Un contrat rapporte la plus haute des trois. Votre contrepartie a regardé l'une des trois et calcule parfaitement~: elle détermine l'espérance du contrat sachant cette carte, vend un lot à votre cours acheteur si celui-ci est supérieur à cette espérance, achète un lot à votre cours vendeur si celui-ci est inférieur, et sinon passe. Cotez un cours acheteur et un cours vendeur écartés de exactement 2, par pas de 0.25. Quelle est la valeur inconditionnelle~? Quelle cotation perd le moins en moyenne, et est-elle centrée sur cette valeur~? Que vaut le contrat après chacune des actions possibles de votre contrepartie~?
\end{question}

\begin{question}[$\star\star\star$ \iqroles{QT, Trader, QR} \iqfirm{teneur de marché, hedge fund} \iqlevel{confirmé}]\label{iq:iv2:making-a-market:31}
Deux dés équilibrés $P$ et $Q$ sont lancés hors de votre vue. Le contrat $S$ rapporte $P + Q$, et le contrat $D$ rapporte $P - Q + 7$. Vous cotez les deux à 6.5 en cours acheteur et 7.5 en cours vendeur. Votre contrepartie a regardé exactement un dé, choisi à pile ou face, et vous ne savez pas lequel. Pour chaque contrat, elle calcule son espérance sachant ce dé, vend un lot si cette espérance est inférieure à votre cours acheteur, achète un lot si elle est supérieure à votre cours vendeur, et sinon passe. Si vous voyez seulement qu'elle a acheté $S$, que vaut chaque contrat~? Que vaut chaque contrat si vous la voyez acheter les deux, ou acheter $S$ et vendre $D$~? Quelle est votre perte moyenne par transaction~?
\end{question}

\begin{question}[$\star\star\star$ \iqroles{QT, Trader, QR} \iqfirm{teneur de marché, hedge fund} \iqlevel{junior, confirmé}]\label{iq:iv2:making-a-market:32}
Un nombre entier de 1 à 10 a été choisi au hasard, et un contrat le rapporte. Vous devez coter, manche après manche, un cours acheteur et un cours vendeur écartés d'exactement 1, tous deux entiers. Votre contrepartie connaît le nombre. À chaque manche, elle vend un lot à votre cours acheteur si le nombre lui est inférieur, achète un lot à votre cours vendeur s'il lui est supérieur, et sinon passe. Comment cotez-vous pour minimiser votre perte totale espérée, quel que soit le nombre de manches, et quelle est cette perte~? Qu'en est-il avec les nombres de 1 à 16~? Combien coûte le fait de commencer à 1 en cours acheteur et 2 en cours vendeur, puis de remonter pas à pas~?
\end{question}

\begin{question}[$\star\star\star\star$ \iqroles{QT, Trader, QR} \iqfirm{teneur de marché, hedge fund, crypto} \iqlevel{confirmé}]\label{iq:iv2:making-a-market:33}
La juste valeur d'un contrat est fixée à 100, et vous cotez toujours 99 en cours acheteur et 101 en cours vendeur, ou retirez un côté de votre cotation. Pendant chacune des 10 manches, un client arrive et achète ou vend un lot à pile ou face, et il échange si ce côté de votre cotation est affiché. Vous commencez à plat. Votre score final est votre profit, la position finale étant évaluée à 100, moins $0.2$ fois le carré de votre position finale. Quand devez-vous retirer un côté~? Donnez, par raisonnement, la position maximale que vous acceptez d'atteindre avec 1, 2 et 3 manches restantes, puis la politique complète. Quel est le meilleur score espéré~? Comparez-le avec le fait de ne jamais retirer un côté, et avec la meilleure limite fixe.
\end{question}

\begin{question}[$\star\star\star\star$ \iqroles{QR, QT} \iqfirm{teneur de marché, hedge fund} \iqlevel{confirmé}]\label{iq:iv2:making-a-market:34}
La vraie valeur $V$ d'une action suit une loi normale de moyenne 100 et d'écart-type 2. Des clients non informés envoient un ordre net $u$, normal de moyenne 0 et d'écart-type 100 actions, indépendant de $V$. Un initié connaît $V$ et, sans voir $u$, envoie un ordre de $x$ actions pour maximiser son profit espéré. Vous ne voyez que le total $y = x + u$, et vous exécutez l'ensemble à un prix qui vous met à l'équilibre sachant $y$. Supposez que l'ordre de l'initié est linéaire en $V$ et que votre prix vaut $100 + \lambda y$. Déterminez $\lambda$, l'ordre de l'initié et son profit espéré, ainsi que la part de la variance de $V$ qui subsiste une fois le prix fixé.
\end{question}

\begin{question}[$\star\star\star\star$ \iqroles{QR, QT} \iqfirm{teneur de marché, crypto} \iqlevel{confirmé}]\label{iq:iv2:making-a-market:35}
Une action se traite à un mid de 100. Vous détenez 2 actions, et en fin de période, le mid suivra une loi normale de moyenne 100 et de variance 4. Votre utilité de la richesse finale est exponentielle, avec une aversion au risque de 0.1, de sorte que l'utilité de la richesse $w$ vaut $-e^{-0.1 w}$. Vous affichez un cours acheteur et un cours vendeur. Exactement un client arrive, voulant acheter ou vendre une action à pile ou face. Un vendeur prend un cours acheteur situé $\delta$ sous 100, et un acheteur prend un cours vendeur situé $\delta$ au-dessus de 100, avec une probabilité $e^{-1.5\delta}$, pour $\delta \ge 0$. Dérivez le cours acheteur et le cours vendeur qui maximisent votre utilité espérée. Quels sont-ils~? Comment se décomposent-ils en un décalage dû à votre position et en une largeur de la fourchette~? Que se passe-t-il quand l'aversion au risque tend vers 0~?
\end{question}
